% Generated draft: requires mathematical and source review.
\documentclass[11pt]{article}
\usepackage[margin=25mm]{geometry}
\usepackage{fontspec}
\setmainfont{Latin Modern Roman}
\usepackage{amsmath,amssymb,mathtools,unicode-math}
\setmathfont{Latin Modern Math}
\usepackage{enumitem,needspace,xurl,hyperref}
\setlength{\parindent}{0pt}
\setlength{\parskip}{5pt}
\setlength{\emergencystretch}{3em}
\newcommand{\N}{\mathbb{N}}
\newcommand{\Z}{\mathbb{Z}}
\newcommand{\Q}{\mathbb{Q}}
\newcommand{\R}{\mathbb{R}}
\newcommand{\C}{\mathbb{C}}
\newcommand{\F}{\mathbb{F}}
\newcommand{\A}{\mathbb{A}}
\newcommand{\PP}{\mathbb{P}}
\newcommand{\E}{\mathbb{E}}
\newcommand{\eps}{\varepsilon}
\newcommand{\dd}{\,\mathrm{d}}
\begin{document}
\section*{UnsolvedMath ID: 800562}
\textbf{Source problem code:} catalogue entry
\subsection{Definitions and mathematical statement}
Let \(P\) be a simple polygon, and let \(V(P)\) be its finite set of vertices. Under the usual closed-region visibility convention, distinct vertices \(u,v\in V(P)\) are visible to each other if the segment \([u,v]\) is contained in the closed polygonal region bounded by \(P\). Its visibility graph is the finite simple undirected graph \(G=(V,E)\), where \(V=V(P)\) and \(\{u,v\}\in E\) exactly when \(u\) and \(v\) are visible. Let \(C\) be the Hamiltonian cycle in \(G\) whose cyclic vertex order follows the boundary of \(P\).

Under the graph-theoretic interpretation of a \emph{hidden set}, a set \(H\subseteq V\) is hidden if its distinct vertices are pairwise nonadjacent:
\[
\forall u,v\in H,\qquad u\ne v\Longrightarrow \{u,v\}\notin E.
\]
Thus a maximum hidden set is a set \(H_{\max}\subseteq V\) satisfying
\[
H_{\max}\text{ is hidden},\qquad
|H_{\max}|=\max\bigl\{|H|:H\subseteq V\text{ is hidden}\bigr\}.
\]
\textbf{Problem.} Given \(G\), promised to be the visibility graph of a simple polygon \(P\), together with its specified boundary Hamiltonian cycle \(C\), determine a maximum hidden set \(H_{\max}\) using the supplied graph and cycle.

\textbf{Unresolved source qualifications.} The supplied statement does not define ``hidden set'' or specify the treatment of boundary contacts in visibility. The formulation above explicitly uses the vertex independent-set interpretation and the closed-region visibility convention; these interpretations require confirmation against the complete source discussion. No coordinates of \(P\) are stated to be supplied. The excerpt specifies no running-time requirement, so no polynomial-time requirement is asserted here.
\subsection{Short English statement}
Given a polygon visibility graph and the cycle recording its boundary order, find a largest set of mutually nonvisible vertices. This interprets a hidden set as an independent set of the given graph. The source excerpt leaves the terminology, geometric visibility convention, and intended computational efficiency unspecified.
\subsection{Sources}
\begin{itemize}
\item Local UnsolvedMath registry, permanent ID \(800562\), source problem code \texttt{catalogue entry}, titled ``Visibility Graphs of Points, Segments and Polygons --- Problem 36.'' The supplied record attributes the problem to Ghosh and Goswami and records the year as 2010.
\item Ghosh and Goswami, \emph{Unsolved Problems in Visibility Graphs of Points, Segments and Polygons} (2010), Problem 36. Supplied source URL: \url{https://arxiv.org/abs/1012.5187}. The record reports access on 2026-10-16 and extraction by \texttt{arxiv\_latex}; these are supplied provenance claims, not independently verified access or source verification.
\item The local record labels the problem \texttt{open} and records an update timestamp of \texttt{2026-10-16T14:00:00Z}. This is supplied registry status, not an independently verified assessment of the literature.
\end{itemize}
\end{document}
